\section{The second split: transmitting a table through the stencil}
\label{inner:section}

At a successful early cut, the suffix can now recover the actual stencil
slot and hence its metadata. That removes the problem of independently
chosen kinds, but an interval crossing the cut is still unreadable. We
handle it using the arbitrary-function form of the shift-table lemma.

The suffix first describes the later computation by a table indexed by
the monoid product of the entire early region. The table has an entry for
every monoid value, not just values realized by prefixes. Composing this
table with the product of the early block values yields a total
function of the tuple. The shift-table lemma supplies a deterministic
affine correction for that function, and one block can then be omitted
for each input state.

\subsection{A total table computed from later letters}

Write $b'_*=b'_{g'}$, the last fixed stencil boundary, and let
$e=[s,b)\in D$.  The homogeneous lift of $F_e$ is the linear map
$\widehat F_e\colon U\to U$, where $U=V\oplus\mathbf F_2$.
If $F_e(v)=vL_e+\gamma_e$, then
\begin{equation}
              \widehat F_e(v,c)=(vL_e+c\gamma_e,c).
 \label{inner:homogeneous-update}
\end{equation}
All marker searches, outer boundaries, and end data occur after
$b'_*$ by Lemma~\ref{schedule:margins}.  Hence the suffix of the
episode starting at $b'_*$, together with its origin and anchors,
determines the following table for every $m'\in M$.

Set $D_b=T_{b'_*,b}$ and $m_b(m')=m'D_b$.  Define an affine map
$H_{e,m'}\colon V\to V$ by the same three cases as $F_e$:
\begin{enumerate}
\item If the actual clock graph is the singleton off-diagonal
      edge $\{(\alpha,\lambda)\}$, put
      \[
      H_{e,m'}(v)=v m_b(m')+Y(\lambda)-X(\alpha)m_b(m').
      \]
\item Otherwise use $u=(q-t)\bmod B$.  If at least one actual
      $f_j^u(t)$ is absent, put $H_{e,m'}(v)=v m_b(m')$.
\item In the remaining case all these searches are present; write their values as
      $l_0<\cdots<l_g$, and put
      \[
      A(m')=m'T_{b'_*,l_0},\qquad C=T_{l_g,b},\qquad
      d_i=T_{l_{i-1},l_i}.
      \]
      With $f_{A(m'),C}$ from \eqref{outer:product-table}, put
      \[
      H_{e,m'}(v)=
      v f_{A(m'),C}(\mathbf d)
      +\beta_{f_{A(m'),C}}(\mathbf d).
      \]
\end{enumerate}
Define
\begin{equation}
 \psi_e\colon M\longrightarrow\operatorname{End}_{\rm lin}(U),
 \qquad \psi_e(m')=\widehat H_{e,m'}.
 \label{inner:total-table}
\end{equation}
The hat denotes the homogeneous lift in
\eqref{inner:homogeneous-update}. For every $m'$, including an impossible
prefix product, associativity in the third case gives
\[
 A(m')d_1\cdots d_g C
   =m'T_{b'_*,l_0}T_{l_0,l_g}T_{l_g,b}
   =m'T_{b'_*,b}=m_b(m').
\]
Thus the affine formula and its linear lift are defined for every entry.
Only after the table has been defined do we substitute the actual early
product, obtaining
\begin{equation}
                    \psi_e(T_{s,b'_*})=\widehat F_e.
 \label{inner:actual-table-entry}
\end{equation}

Every letter-dependent quantity here is read after $b'_*$: marker
positions, anchor data, inter-boundary products, and products from
$b'_*$ to later positions. The hypothetical value $m'$ is simply a
finite algebraic input. The calculation does not search for omitted
prefix letters, and it does not query graph languages for sequences of
episodes.

Apply exactly the same formulas under a proposed suffix origin once its
end rule and all necessary data are available. If a required interval is
missing or its endpoints are in the wrong order, reject the entire suffix
test. Do not discard that origin, and do not discard individual values of
$m'$ from the table. This separates the availability of actual suffix
letters from the total algebraic domain of $\psi_e$. At the true origin
all required intervals are available by the margin lemma.

\subsection{The affine update on the exceptional domain}

For $1\le i\le g'$ put
\[
                         d'_i=T_{b'_{i-1},b'_i}.
\]
For any table $\psi\colon M\to\operatorname{End}_{\rm lin}(U)$,
define
\begin{equation}
 f'_\psi(a_1,\ldots,a_{g'})=\psi(a_1\cdots a_{g'}).
 \label{inner:product-table}
\end{equation}
Use the fixed shift table $\beta_{f'_\psi}$ supplied by
Lemma~\ref{alg:shift-table} with state space $U$ and dimension
$g'$.  On every $e\in D$ define
\begin{equation}
 G_e(x)=x f'_{\psi_e}(\mathbf d')+
                       \beta_{f'_{\psi_e}}(\mathbf d'),
 \qquad x\in U.
 \label{inner:affine-update}
\end{equation}
This is a deterministic total affine map on $U$.  Its linear
part is $\widehat F_e$, because
$d'_1\cdots d'_{g'}=T_{s,b'_*}$ and
\eqref{inner:actual-table-entry} applies.

The kinds in the stencil were chosen to contain all tuples
\[
          (j,x,y,\mathbf d'_{-j},\psi),
\]
where $1\le j\le g'$, $x,y\in U$,
$\mathbf d'_{-j}\in M^{g'-1}$, and
$\psi\colon M\to\operatorname{End}_{\rm lin}(U)$ is arbitrary.
The $\mu$ copies of a kind have distinct cut offsets but the same
fields.  A cut of block $j$ lies strictly inside
$(b'_{j-1},b'_j)$, so all listed block products except the omitted
$j$th product can be checked on one side or the other.

For $x,y\in U$ define inner prefix and suffix languages
$\widetilde P_x,\widetilde Q_y$ as follows.
A prefix chooses a slot $p\in\mathcal P$ and has exactly length
$p$, so its endpoint is $k=s+p$.  If the slot's kind is
$(j,x_p,y_p,\mathbf a_{-j},\psi)$, it requires $x_p=x$ and
checks all listed earlier block products $a_i=d'_i$ for $i<j$.
For $a\in M$, write $\mathbf a[j\leftarrow a]$ for the tuple
formed by inserting $a$ into the listed entries at position $j$.
The prefix also checks the finite certificate
\begin{equation}
 x_p f'_\psi(\mathbf a[j\leftarrow a])
 +\beta_{f'_\psi}(\mathbf a[j\leftarrow a])
 =y_p\qquad(a\in M).
 \label{inner:certificate}
\end{equation}
The prefix does not try to check the later entries or the table
$\psi$ against letters it cannot see.

A suffix beginning at $k$ performs the complete decoder of
Lemma~\ref{decode:decoder}, including its reference end guard
when called for and its final decoded-origin end guard.
Let $p$ be its decoded slot and use its kind
$(j,x_p,y_p,\mathbf a_{-j},\psi)$.  At the decoded origin the
suffix computes the actual searches, anchors, and all entries
of \eqref{inner:total-table}.  It requires that resulting table
equal the slot's $\psi$, checks $a_i=d'_i$ for $i>j$, and requires
$y_p=y$.  The comparison of tables is equality on every input
$m'\in M$.  Each required interval starts on or after the
suffix's beginning: in particular, the displacement of $b'_*$
from $k$ is the fixed positive offset $(b'_*-s)-p$, since every
stencil slot precedes the last stencil boundary.  The margin lemma
places all the data for the table after that boundary.

\begin{proposition}[Second split]\label{inner:split}
For the updates $G_e$ on $D$, the languages
$\widetilde P_x,\widetilde Q_y$ satisfy:
\begin{enumerate}
\item on a word in $D^*$, every match of
      $\widetilde P_x\widetilde Q_y$ from its beginning consumes
      exactly the first episode $e$ and has $G_e(x)=y$;
\item every $e\in D$ admits such a split for every input $x\in U$.
\end{enumerate}
Thus this split has empty skip domain.
\end{proposition}

\begin{proof}
For soundness, let an arbitrary successful pair begin a word of $D^*$.
The prefix selects a scheduled stencil length, so its endpoint $k=s+p$
lies inside the first episode. The decoder lemma forces the suffix to
select that very slot and to end at the actual episode endpoint. Thus
both factors use one kind, although this was not assumed when defining
their languages.

The prefix verifies every listed early block product, and the suffix
verifies every listed later one. It also compares its computed table
with the kind's table on \emph{all} $m'\in M$. Hence the listed products
are the actual $d'_i$ for $i\ne j$, and the table is $\psi_e$.
The actual product $d'_j$ of the interval crossing the cut is a member
of $M$, so it is covered by the universal certificate
\eqref{inner:certificate}. Substituting it proves $G_e(x)=y$.
The positive stencil length rules out a pair on the empty word.

For completeness, fix an episode $e\in D$ and any input $x\in U$.
Apply Lemma~\ref{alg:shift-table} to the total function
$f'_{\psi_e}$ at its actual tuple $\mathbf d'$. A coordinate $j$
exists on which the shifted output is constant, equal to $G_e(x)$.
Take the kind with that coordinate, the actual other products, the entire
actual table $\psi_e$, and input and output fields $x,G_e(x)$. It belongs
to the roster because all total tables were included when kinds were
chosen. Its finite certificate holds for every possible missing value.
The decoder lemma leaves at least one usable copy of this kind. At that
copy all product checks concern actual readable intervals, its table
comparison is exact, and its end checks pass. The split therefore exists
for the chosen $x$. Repeating this reasoning for each input proves the
all-input claim; neither the omitted coordinate nor the chosen copy is
required to be the same for different inputs.
\end{proof}

We have now obtained both split constructions with their full
all-input guarantees.  The inner updates $G_e$ have linear
parts $\widehat F_e$, and the outer updates $F_e$ have linear
parts equal to the original monoid action.  The remaining task
is to express each domain and each individual prefix or suffix
test with the claimed height over the original alphabet, and
then apply affine recovery and the split identity in this
order.
